%Yap master Lili \cite{Hinton}
We constantly have something on our minds – the neurons in our brains are turbulently exchanging information. This information is in the form of electrical signals, and scientists have long ago discovered a way of non-invasively recording this activity – electroencephalogram, referred to as EEG (\cite{ROSSINI2025132}). EEG is a method of measuring continuous patterns of electrical activity in underlying regions of the brain (\cite{eeg_measuring_method}). To record this activity, electrodes are positioned on different areas of the scalp. 
\vspace{0.5em}\\
Since EEG is significantly advantageous in multiple ways, it has acquired a wide use across multiple fields, some of which being psychophysiology, cognitive, and neuroscience. The analysis of EEG consists of combining mathematical methods and computer technology, whose synergy has led to researchers gaining a deeper understanding of the inner workings of the brain. Typically, EEG analysis methods are categorized as falling into the time, frequency, time-frequency, or non-linear domains. In addition, with the development of deep learning, a new category of analysis has emerged – namely, EEG research and analysis via deep artificial neural networks. 
\vspace{0.5em}\\
There are many kinds of neural networks, thus, many ways of analysing EEG signals. One of them, which has acquired a wide use among researchers and scientists, is the graph neural network (GNN). In general, GNNs have proved to be well-suited for dealing with EEG data, as they are able to model complex structures, which are inherent in brain connectivity (\cite{Song}). A specific type of a GNN, which was shown to be a great tool for EEG analysis and research, is the dynamic graph convolutional neural network (DGCNN). It can adaptively learn the innate relationships of EEG channels, allowing the algorithm to produce a more accurate graph representation of the EEG data (\cite{Song}). Graph representations, in turn, offer an intuitive way of modelling brain activity. In the case of working with EEG -- representing electrodes as nodes and their interactions as edges allows for the capturing of patterns of functional connectivity between different brain regions. This structure-oriented approach aligns well with the brain’s inherently networked organization and opens the door to leveraging a wide range of tools from network analysis. Via them, we may manage to reveal interpretable insights about the activity from and between functionally different brain areas. 
\vspace{0.5em}\\
However, for such insights to be meaningful, it is desirable, that the learned graphs are stable, meaning that their structures are consistent. If repeated training runs on the same data produce significantly different graph structures, it becomes difficult to draw conclusions and trust the model’s graph constructing ability. This instability can significantly undermine the reliability of the model and the produced results.
\vspace{0.5em}\\
Despite these challenges, the potential of EEG analysis, carried out with the aid of artificial neural networks, remains compelling. Our research deals with investigating the consistency of graph structures learned by GNNs. Specifically, we explore methods for measuring representational similarity and identifying persistent patterns across training runs, aiming to better understand how reliable these learned graphs truly are. 